\section{对象、原生组织与服务边界}

本例属于 3D FeNOR／vertical FeFET 技术族，主对象为 \sid{FENOR-02} 的 \textbf{2026 3D vertical AND FeFET}。采用 6 nm SL-HZO、4 nm O-poor IGO 的二状态切换，并将 O-0.1s 工艺的扰动抑制作为相容参考条件。器件、局部阵列和仿真为选型提供依据，完整 INT8 服务由明确的 CMOS 资源和调度构成。

所选局部组织存放 $W[N,K]$，$K=N=128$、$b_S=b_R=1$ Byte/元素；每向量输入 $B_S=128$ Byte，完整 resident 容量为 16384 Byte。每个 INT8 权重使用八个二状态 cell，最高位在数字重构中取 $-2^7$。输出容器为 $16+\lceil\log_2K\rceil=23$ bit，每次交付 128 个结果。该尺寸是本例的组织选择，其他技术可以采用不同原生规模。

\begin{table}[htbp]\centering\small
\begin{tabularx}{\textwidth}{@{}p{24mm}X@{}}\toprule
来源与定位 & 证据及采用方式\\\midrule
\sid{FENOR-02} p.2 Figs.3--5 & SL+O-poor 的 $\pm2$ V、20 ns 切换；写后读延迟 RAWD$<100$ ns；$-0.5$ V、20 ns 读耐久应力。脉冲平台、观察预算和完整本地读周期分别处理。\\
\sid{FENOR-02} p.3 Figs.9--11 & $V_w/3$ 邻层偏置，最高 $2V_w/3$ 半选；$4\times16\times16$ cell 图案验证；层数扩展读延迟来自 TCAD 提取 RC 后的 SPICE。作为选通、局部规模和读侧量级依据。\\
\sid{FENOR-01} p.3 Fig.12 & 2025 NOR-type 的横向多 cell 权重、SA、寄存器和加法器。用于二进制读码加数字归约的机制桥接；其 50 ns 写脉冲与约 100 nA 读电流不进入主结果。\\
\sid{FENOR-04/06} & 分别为 ZnO/MFMIS 计算阵列和不同栅堆栈的速度／可靠性比较。保留其结构、偏置与终点差异，不将多级脉冲或其他栅堆栈速度拼入本例。\\\bottomrule
\end{tabularx}
\caption{关键证据入口。页码按本地 PDF 页序，完整原值及条件见证据笔记。}
\end{table}

\sid{FENOR-02} 制造尺寸为 $4\times32\times32$，Fig.10 图案验证为 $4\times16\times16$；本文布局是基于其选通关系的局部参考宏。SL+O-poor 切换与 O-0.1s 扰动结果的组合条件、4096 路感测和写驱动仍需电路及工艺验证。\sid{FENOR-01/02/06} 为同团队证据。两路边界包括该宏的阵列、驱动、感测、有限保持、数字归约和结果寄存器；上游传输不包含在本地服务间隔中。

\clearpage
\section{逻辑映射、有限保持与完整求值}

\subsection{层承载容量，横向读通路提供并行}
构造 32 条横向行读通路，每条的四层存放四个输入索引；共享 BL/SL 的层每次只选一个。每条行通路包含 64 个局部条带，即八个输出组乘八个位平面。一个条带有四层、每层 16 个 BL/SL 列对，容量 64 cell。映射为
\[
W_{o,i}[b]\longmapsto
(i\bmod32,\;\lfloor i/32\rfloor,\;8\lfloor o/16\rfloor+b,\;o\bmod16).
\]
坐标依次为横向行通路、层、条带和列；总物理容量为
\[
32\times64\times4\times16=131072\ \text{binary cell}
 =16384\ \text{有效逻辑 Byte}.
\]
没有信号增强复制，也没有将位平面作为独立服务副本。16 列局部条带与已测阵列尺度相接；32 行通路是实现声明活动宽度所需的真实读资源。

每条横向行通路配 128 个二进制感测节点，共 $32\times16\times8=4096$ 个 SA；八组输出复用这些 SA。未选层通过局部关断偏置隔离。一批得到 $32\times16\times8=4096$ bit 权重，捕获进单个 4096-bit tile 寄存器。随后连续八个输入位使用这份读码，读码释放后才读取下一 tile；不缓存整份权重，也不边计算边覆盖单 bank。

16 条数字通道每拍完成各自 32 项的位乘法、归约和位权累加。输入寄存器为 $128\times8=1024$ bit；16 个 23-bit 部分和寄存器跨行组保持，完整输出寄存器为 $128\times23=2944$ bit。各输出组完成全部行组后提交，所有输出就绪才完成一个向量服务。输出不增加 $Q_S$，其计算和提交时间仍计入服务。

\subsection{读预算和阶段计数}
完整局部读响应 $t_m=20/40/80$ ns 覆盖选择、读偏置／RC、二进制感测和读后恢复。20 ns 读应力和 Fig.11 的 ns 量级 RC 仿真为量级锚点，三值为固定组织上的工程服务条件；并非完整宏的实测快慢边界。新增数字捕获每物理读批另付一 $T_D$，避免将未被原周期覆盖的锁存隐入 $t_m$。

共同数字拍为 $T_D=2/5/10$ ns。四个行组、八个输出组和八个输入位给出
\[
N_{\rm rd}=4\times8=32,\quad N_{\rm cap}=32,\quad N_D=32\times8=256,
\]
\[
\boxed{\Delta_S=32t_m+32T_D+256T_D+2T_D.}
\]
最后两拍为整向量输入捕获和输出提交。无跨 tile 流水隐藏时间；典型为 $1280+160+1280+10=2730$ ns。物理读与数字归约各约占 47\%，有限保持使二者共同决定输入服务能力。此数字路径无 ADC，也不附加 ACIM 的输入驱动／转换时隙。

\clearpage
\section{128-cell 更新资源与完整终点}

外部更新采用单域 128-bit 编码接口及 128-bit 目标数据寄存器。一笔对齐事务对应一个输入索引和连续 16 个输出，$B_R=16$ Byte。八个位平面条带同时驱动，共 128 个二元目标；该并行度由选通、驱动和负载预算支持，接口位宽本身不提供物理并行。一次完整矩阵装载为 $128\times(128/16)=1024$ 笔事务。

\subsection{两相掩码与驱动条件}
令 $s=\pm1$、$V_w=2$ V。沿 \sid{FENOR-02} Fig.9 的偏置关系，本相目标列与其余抑制列分别设置为：

\begin{table}[htbp]\centering\small
\begin{tabular}{@{}lccc@{}}\toprule
cell 类别 & WL & BL=SL & 栅沟道差\\\midrule
本层、本相目标 & $2sV_w/3$ & $-sV_w/3$ & $sV_w$\\
本层、抑制列 & $2sV_w/3$ & $sV_w/3$ & $sV_w/3$\\
其他层、目标列 & $sV_w/3$ & $-sV_w/3$ & $2sV_w/3$\\
其他层、抑制列 & $sV_w/3$ & $sV_w/3$ & $0$\\\bottomrule
\end{tabular}
\caption{独立局部条带内的多列掩码。两相极性相反，邻层 WL 差均为 $V_w/3$。}
\end{table}

正相将目标 bit 1 写入，负相将目标 bit 0 写入，每极性平台 20 ns。先写成 1 的列在负相只承受半选场。普通混合目标字按两个非空掩码相计时，不假定每 cell 都经过两次全幅切换，也不需要独立块擦除。Fig.9 的最高 $2V_w/3$ 扰动结果支持这种场分类；同时多目标的波形、扇出及供电为参考电路条件。

八条带需要 $128$ BL、$128$ SL 和 $32$ WL，共 288 个独立驱动节点。选每节点有效负载上限 100 fF、最坏摆幅 $8/3$ V，则边沿 $t_e=5/10/20$ ns 分别需要
\[
I_{\rm node}\ge C\Delta V/t_e=53.3/26.7/13.3\ \mu\mathrm A,
\]
总瞬态需求为 15.36/7.68/3.84 mA。三情景保持同一 288 节点、每节点至少 60 $\mu$A、总额定能力至少 17.28 mA 的本地资源；典型 7.68 mA 是 10 ns 边沿的需求，不是供给上限。100 fF 为局部短条带的待验证设计上限；Fig.11 横向 120/170 nm 节距和实测 40 nm 竖直节距支持局部负载组织，但未给出本宏提取电容。本地偏置轨已就绪，电源冷启动不在本地边界内。

\subsection{观察、捕获、比较只计一次}
最后一个极化平台结束后取 $g=100$ ns 观察窗口，包含最终写偏置退回。其依据是 RAWD$<100$ ns 所支持的保守观察选择，而非材料固有必需等待。窗口之后执行完整二进制读 $t_m$，另以一拍将 128 bit 捕获进既有 tile bank，再由 16 路比较器用八拍终验。更新期间 streaming 停止，终验复用单 bank，不另加 shadow。

顺序为“正相建立及平台、相间退回、负相建立及平台、含最终退回的观察、读回、捕获、比较”。共同接口／完成控制共两拍，因此
\[
\boxed{\Delta_R=2T_D+3t_e+2(20)+\max(g,t_e)+t_m+T_D+8T_D.}
\]
典型分项为 $10+30+40+100+40+5+40=265$ ns。终验失败作为异常，未假定成功率或额外重试数；若应用要求重试，需按实际服务追加。非对齐小更新按完成的有效字节计算，不自动获得完整 16-Byte 分子。

\clearpage
\section{典型能力、有限情景与两表接口}

采用逻辑 payload 与对应服务，$\rho=B_S/\Delta_S$、$\tau=B_R/\Delta_R$。典型得到
\[
\rho=46.89\ \mathrm{MB/s},\quad\tau=60.38\ \mathrm{MB/s},\quad
\RI^*=\frac{128}{16}\frac{265}{2730}=0.7766.
\]
% Generated by scripts/check_fenor.py --emit.
\begin{table}[htbp]\centering\small
\begin{tabular}{@{}lrrrrr@{}}\toprule
情景 & $\Delta_S$ (ns) & $\Delta_R$ (ns) & $\rho$ (MB/s) & $\tau$ (MB/s) & $\RI^*$\\\midrule
乐观 & 1220 & 197 & 105 & 81.2 & 1.29\\
典型 & 2730 & 265 & 46.9 & 60.4 & 0.777\\
悲观 & 5460 & 390 & 23.4 & 41 & 0.571\\
16-cell 资源对照 & 2730 & 1770 & 46.9 & 9.04 & 5.19\\
\bottomrule\end{tabular}
\caption{所选原生 $128\times128$ INT8 配置；主情景固定128-cell写资源、4096-bit单tile保持与100 ns观察窗口。最后一行独立减少写驱动，MB=$10^6$ Byte。}
\label{10_fenor_3d:tab:result}\end{table}


主三情景使用相同层数、SA、保持、归约、写驱动额定能力与 $g=100$ ns，仅配对本地读响应、偏置边沿及共同数字服务条件。典型值取合理中间服务条件；快慢值是有限工程情景，不能解释为概率区间或圈内所有参数组合均可实现。短情景尤其依赖安装驱动在实际负载下支持 5 ns 边沿，以及 SA 与数字通道满足相应时间。

\subsection{完整装载与平衡复用次数}
1024 笔完整对齐事务覆盖本矩阵，每个有效权重恰更新一次：
\[
Q_R=16384\ \mathrm{Byte},\quad T_R=1024\times265=271360\ns.
\]
对一次装载后 $U$ 个完整向量，$Q_S=128U$ Byte，$\RI=U/128$；匹配两路的平衡次数为
\[
U^*=\frac{T_R}{\Delta_S}=99.40=128\RI^*.
\]
$U^*$ 使用完整矩阵装载与向量服务间隔，不使用小事务的 $\Delta_R/\Delta_S$ 替代。该配置顺序执行，所以服务间隔等于完整本地求值时间；同时读写互斥，两路峰值不保证同时达到。更大算子应按输入共享、分时和更新域重新聚合 $T_R,\Delta_S$，不将此 $U^*$ 无条件转移到其他输出维度或复合阶段。

本例有效容量恰为 16 KiB，因此每 16 KiB 平均更新成本与完整装载时间数值相同；一般换算为 $T_R\,16384/(KNb_R)$，不把平均成本当任意请求延迟。

\subsection{资源和观察窗口的独立对照}
若只安装单条带写资源，即 36 个节点、每节点同样的额定能力，一次只更新 16 cell，八个位平面顺序完成。每批保留完整两相、观察、读回、一拍捕获及一拍比较：
\[
\Delta_{R,16}=10+8(30+40+100+40+5+5)=1770\ns.
\]
$\rho$ 不变，$\tau=9.040$ MB/s、$\RI^*=5.187$。此差异来自真实写驱动和调度，不是层数或材料速度改变；它单列为资源对照。

观察窗口也是重要工程判断。固定典型所有资源与服务条件时，$\Delta_R(g)=165+g$ ns，$g\ge10$ ns。取额外预留 $g=150$ ns 得 $\Delta_R=315$ ns、$\tau=50.79$ MB/s、$\RI^*=0.9231$；150 ns 不属于普通悲观点，也不是实测上界。典型观察占写服务约 38\%，终态读回、捕获与比较占约 32\%；20 ns 极化平台本身不能代表完整更新能力。

\clearpage
\section{维护、证据边界与复算入口}

所选非易失二进制状态在本地服务窗口无周期刷新或破坏性读恢复，机器结果因此保存相同的维护前和长期局部有效能力。原文写耐久 $>10^{12}$、读应力耐久 $>10^{12}$、半选耐受 $10^9$ 以及约 3000 s 保持测量／十年外推分别对应自己的条件。长期半选累计、阈值漂移和寿命需另行评估，不将耐久数目折成瞬时 $\tau$。

本例位置主要由三项决定：4096 SA 加有限 tile 保持将物理读次数压到 32，同时仍完成 256 轮数字归约；八条带实际驱动使八个位平面能同批更新；观察与完整终验占据写侧的大部分时间。不同层只是 resident 容量，较小 $\RI^*$ 也不能单独证明某种动态更新 workload 的性能优势。

\subsection{原始文献定位}
以下按本地 PDF 页序；VLSI 两篇为三页本地稿件，无文件内版本／日期标识。
\begin{enumerate}
\item \sid{FENOR-02}，Y. Zhou 等，\emph{3D Vertical FeFET Array with Record Endurance, Fast Writing, Disturb Immunity, and Kb-scale Verification}，VLSI 2026。p.1 结构／工艺，p.2 Figs.2--5 尺寸、脉冲、RAWD 和耐久，p.3 Figs.7--10 半选与图案，Fig.11 TCAD/SPICE。DOI：\href{https://doi.org/10.1109/VLSITechnologyandCir65830.2026.11577291}{10.1109/VLSITechnologyandCir65830.2026.11577291}。
\item \sid{FENOR-01}，Y. Zhou 等，\emph{3D NOR-type FeFETs with Record Endurance, Fast Erase and Immediate Read-After-Write}，VLSI 2025。p.2 Figs.2--5，p.3 Fig.12 数字计算机制。DOI：\href{https://doi.org/10.23919/VLSITechnologyandCir65189.2025.11074820}{10.23919/VLSITechnologyandCir65189.2025.11074820}。
\item \sid{FENOR-04}，Y. Feng 等，\emph{Efficient Large Scale Neural Network Acceleration With 3-D FeNOR-Based Computing-in-Memory Design}，IEEE TED 72(5)，2025。p.4 Figs.7--8，p.5 Figs.9--11；pp.5--6 为模型。DOI：\href{https://doi.org/10.1109/TED.2025.3554164}{10.1109/TED.2025.3554164}。
\item \sid{FENOR-06}，Y. Zhou、R. Huang、K. Tang，\emph{Gate Stack Engineering of 3D Oxide Channel FeNOR Memory}，EDTM 2026。p.2 Sec.II-C 的 90\% 终点，p.3 Figs.8--10 不同栅堆栈速度与可靠性。DOI：\href{https://doi.org/10.1109/EDTM65772.2026.11498024}{10.1109/EDTM65772.2026.11498024}。
\end{enumerate}

\subsection{机器数据与外部审阅条件}
输入文件 \texttt{data/inputs.json} 保存原生维度、物理资源、阶段覆盖、原值与工程选择。结果文件 \texttt{data/results.json} 保存未取整结果、原生容量、完整装载映射接口和来源／共享 API 哈希。\texttt{scripts/check\_fenor.py} 默认只读，\texttt{--emit} 更新结果和表格。检查包括全 cell 唯一映射、两极性电场分类、单 tile 生命周期、驱动额定能力、完整读写账本、终验捕获、矩阵覆盖和 $U^*$ 关系。

需要外部器件／电路审阅的主要判断是：4096 SA 和 32 项数字归约的负载与时序；SL+O-poor 和 O-0.1s 的工艺组合；288 节点在 100 fF 上限下的多目标波形、半选与供电；100 ns 观察和一次终验的充分性。这些条件通过资源与服务预算公开表达，内部算术检查不替代其物理验证。
